\chapter{Bayesian Model and Diagnostic Details}
\label{app:priors}\label{app:poisson}\label{app:likelihood}\label{app:diagnostics}

\section{Coefficient Priors}
The three Bayesian models share their survival likelihood, predictor matrices, baseline intervals and validation folds. The Normal prior is $\beta_j\sim\mathcal{N}(0,1)$. The horseshoe arm uses global and local half-$t_2$ scales and a regularising slab. The continuous spike-and-slab arm mixes spike and slab variances using $\gamma_j\sim\operatorname{Beta}(1,1)$. Equations~\eqref{eq:normal-prior}--\eqref{eq:spike-slab-prior} give the complete specifications. Continuous mixing weights and local shrinkage summaries describe their respective priors; posterior coefficient and hazard-ratio summaries provide the shared interpretation across families.

\section{Interval Exposure and Prediction}
Let $a$ be the median observed event time in the current training partition. Follow-up is scaled as $u=t/a$. Six baseline intervals are formed from training-event quantiles, with an initial boundary of zero and a final boundary beyond the maximum scaled training follow-up. Equally spaced boundaries are used when tied quantiles produce insufficient distinct boundaries.

Patient $i$ contributes exposure
\[
e_{im}=\max\{0,\min(Y_i/a,b_m)-b_{m-1}\}
\]
and event count
\[
d_{im}=\delta_i\mathbb{I}\{b_{m-1}<Y_i/a\leq b_m\}
\]
in interval $m$. Positive-exposure rows enter a Poisson model with mean $e_{im}\lambda_m\exp(\boldsymbol{x}_i^{\mathsf T}\boldsymbol{\beta})$. This representation yields the full piecewise-exponential survival likelihood up to parameter-independent factors. An event on a boundary belongs to the interval ending at that boundary.

Prediction integrates the interval hazards for each posterior draw and averages the resulting survival probabilities. In original time units, the hazard includes the scaling factor $1/a$. The last interval hazard continues beyond the final boundary. The four- and ten-interval models provide the baseline-resolution sensitivity comparisons.

\section{Sampling Diagnostics}
Each initial fit uses four chains, with 2,000 warm-up iterations and 2,000 retained draws per chain. Effective sample sizes account for autocorrelation within these draws. The diagnostic criteria and two specified refits are described in Chapter~\ref{chap:methods}. Acceptance requires every component of the diagnostic criterion to be satisfied.

The primary Bayesian analysis contains 205 accepted outcomes out of 225 model--fold evaluations. The robustness analysis contains 770 accepted outcomes out of 825 evaluations. Table~\ref{tab:gen-robustness-diagnostics} details diagnostic outcomes by cohort and prior family. Predictive and paired-effect calculations use accepted fits.

% Generated from hash-verified saved evidence. Do not edit by hand.
\begin{table}[htbp]
\centering
\caption{Diagnostic outcomes in the robustness analyses.}
\label{tab:gen-robustness-diagnostics}
\small
\setlength{\tabcolsep}{3pt}
\begin{adjustbox}{max width=\textwidth}
\begin{tabular}{llrrr}
\toprule
Cohort & Prior & Evaluations & Accepted & Other outcomes \\
\midrule
GBSG2 & Normal & 85 & 85 & 0 \\
GBSG2 & Horseshoe & 85 & 45 & 40 \\
GBSG2 & Spike-and-slab & 85 & 85 & 0 \\
WHAS500 & Normal & 90 & 90 & 0 \\
WHAS500 & Horseshoe & 90 & 87 & 3 \\
WHAS500 & Spike-and-slab & 90 & 90 & 0 \\
METABRIC & Normal & 100 & 100 & 0 \\
METABRIC & Horseshoe & 100 & 88 & 12 \\
METABRIC & Spike-and-slab & 100 & 100 & 0 \\
\bottomrule
\end{tabular}
\end{adjustbox}
\par\smallskip
\begin{minipage}{\textwidth}\footnotesize Counts describe all 825 robustness conditions. Paired effects use diagnostically accepted fits.\end{minipage}
\end{table}

